\documentclass[preprint,12pt,authoryear]{elsarticle}

\usepackage{amssymb}
\usepackage{amsmath}
\usepackage{mathrsfs}
\usepackage[utf8]{inputenc}
\usepackage[T1]{fontenc}
\usepackage{CJKutf8}
\usepackage{booktabs}
\usepackage{threeparttable}
\usepackage{graphicx}
\usepackage{subcaption}
\usepackage{rotating}
\usepackage{tabularx}
\usepackage{array}
\usepackage{float}
\usepackage{xcolor}



\journal{Financial Research Letter}

\begin{document}

\begin{frontmatter}

\title{Global Geopolitical Risk, Ambiguity, and Emerging Market Returns: Evidence from China}

\author[inst1]{He Ni}
\affiliation[inst1]{organization={Tailong School of Finance, Zhejiang Gongshang University},
addressline={18 Xuezheng Str.},
            city={Hangzhou},
            postcode={310018},
            state={Zhejiang},
            country={China}}

\author[inst2]{Muhammad Ansar Majeed}
\affiliation[inst2]{organization={Department of Accounting, Finance and Economics, KEDGE Business School},
addressline={Domaine de Luminy},
            city={Rue Antoine Bourdelle},
            postcode={13009},
            state={Marseille},
            country={France}}

\author[inst3]{Xingjian Guo}
\affiliation[inst3]{organization={School of Finance, Zhejiang Gongshang University},
            addressline={18 Xuezheng Str.},
            city={Hangzhou},
            postcode={310018},
            state={Zhejiang},
            country={China}}

\begin{abstract}
We study how global geopolitical risk (GPR) affects Chinese equity returns through two distinct uncertainty channels: risk (volatility) and ambiguity (Knightian uncertainty). Ambiguity is conceptually and empirically independent of volatility; while less commonly used than volatility, it adds incremental explanatory power for returns. We construct daily ambiguity measures under an adaptive model uncertainty framework and show that GPR significantly raises financial ambiguity, and higher ambiguity is associated with lower returns. When ambiguity is included alongside traditional risk, models explain more variation in returns and the role of volatility-based risk is often reduced. Our integrated analysis—combining time-series tests with mediation and robustness checks—provides a comprehensive account of how geopolitical risk transmits to asset prices not only through risk but also through ambiguity.
\end{abstract}

\begin{keyword}
Geopolitical Risk \sep Ambiguity Aversion \sep Asset Pricing \sep Chinese Capital Market 
\end{keyword}

\end{frontmatter}

\section{Introduction}

Geopolitical risk (GPR) is a pervasive driver of uncertainty in emerging markets. We argue that uncertainty operates through two distinct channels: risk (volatility) and ambiguity (Knightian uncertainty). Ambiguity is conceptually and empirically independent of volatility \citep{Knight1921, Ellsberg1961}; in asset pricing, models of ambiguity aversion \citep{Gilboa1989, Klibanoff2005} imply that higher ambiguity lowers valuations even when volatility is unchanged.

Our goal is to provide a comprehensive account of how GPR affects returns through both channels. We construct our daily ambiguity measures using a novel, simplified cross-entropy-based approach, which we argue is empirically more robust than traditional proxies like the variance of variance. This methodological contribution, detailed in Appendix A, is the foundation for our empirical tests.

Our empirical contributions are threefold. First, we show that our new ambiguity measure is negatively priced in both time-series (CSI 300) and cross-sectional (constituents) tests, providing incremental explanatory power beyond volatility. Second, we establish that GPR significantly increases financial ambiguity, identifying it as a key transmission channel from geopolitics to asset prices \citep{Caldara2022}. Third, in an integrated specification, the inclusion of ambiguity improves model fit and often reduces the explanatory power of volatility-based risk. We further explore heterogeneity in the pricing of ambiguity by testing for moderation effects related to state-owned enterprise (SOE) status.

Methodologically, our empirical strategy is designed to rigorously test the links between geopolitical risk, ambiguity, and returns. We employ time-series regressions to establish the direct impact of GPR on our daily ambiguity measures and to confirm that ambiguity commands a negative risk premium over time. To formally test the transmission channel, we utilize mediation analysis to quantify the extent to which ambiguity explains the effect of GPR on returns. Furthermore, we conduct Fama-MacBeth cross-sectional regressions to verify that ambiguity is a systematically priced factor across individual stocks. A key part of our analysis is exploring heterogeneity through moderation tests. We investigate whether the pricing of ambiguity differs for state-owned enterprises (SOEs) versus non-SOEs. This tests the hypothesis that implicit government guarantees may buffer SOEs from the adverse effects of ambiguity, leading to a weaker negative relationship between ambiguity and returns for these firms. The robustness of our findings is verified through a comprehensive set of tests. To address potential endogeneity concerns, such as reverse causality from returns to ambiguity, we employ an instrumental variable (IV) approach. We also conduct event studies around major geopolitical episodes (e.g., escalations in the US-China trade conflict) to assess the immediate impact of GPR shocks on ambiguity in a quasi-natural experimental setting. Further checks include using alternative specifications for our ambiguity measure, controlling for a wider array of firm-level characteristics in the cross-section, and examining the stability of our results across different sub-periods.

The remainder of the paper is organized as follows. Section \ref{sec:literature} reviews the literature and develops our hypotheses. Section \ref{sec:data} first details the data and empirical methodology for exploring global geopolitical risk’s transmission mechanism in Chinese equity markets. It then presents the main results and corresponding robustness checks to confirm these findings’ reliability., and Section \ref{sec:conclusion} concludes.

\section{Literature Review and Hypotheses Development}
\label{sec:literature}

Our research is situated at the intersection of two burgeoning fields: the asset pricing implications of ambiguity and the impact of geopolitical risk (GPR) on financial markets. We build a bridge between them by proposing that ambiguity is a primary, yet under-appreciated, channel through which geopolitical tensions transmit to equity returns.

\subsection{Geopolitical Risk and Financial Markets}

The systematic study of geopolitical risk's economic impact has been significantly advanced by the development of quantitative measures, most notably the GPR index of \citet{Caldara2022}. This index, derived from news text, has catalyzed a wave of research demonstrating that elevated GPR is detrimental to financial markets. The consensus finding is that GPR shocks lead to lower stock returns, higher volatility, and a flight to safe-haven assets \citep{Caldara2022}. The effects are often amplified in emerging markets, which are more sensitive to shifts in global capital flows and investor sentiment.

However, the literature has predominantly focused on risk, proxied by volatility, as the main transmission mechanism. While studies have linked political uncertainty to reduced investment \citep{Pastor2013} and economic policy uncertainty to depressed economic activity \citep{Baker2016}, the specific role of ambiguity—or Knightian uncertainty, where probabilities are unknown—remains largely unexplored. This represents a significant gap. Geopolitical events are, by their nature, fraught with ambiguity; they create situations where past data is a poor guide to the future, and the range of possible outcomes and their likelihoods are difficult to assess. We argue that failing to account for ambiguity leaves our understanding of the GPR-return nexus incomplete.

\subsection{Ambiguity Aversion and Asset Pricing}

The theoretical distinction between risk (known probabilities) and ambiguity (unknown probabilities) dates back to \citet{Knight1921} and was famously demonstrated by \citet{Ellsberg1961}. Modern asset pricing theory has incorporated this distinction through models of ambiguity aversion, such as the multiple-priors framework of \citet{Gilboa1989} and the smooth ambiguity model of \citet{Klibanoff2005}. These models predict that ambiguity-averse investors demand a premium for holding assets with uncertain payoff distributions, which depresses asset prices independent of conventional risk.

Empirically testing these theories requires a credible measure of ambiguity. The literature has proposed several proxies, such as the dispersion of professional forecasts \citep{Brenner2018, Ulrich2013} or the variance of variance. However, these proxies can be indirect and may not be available at a high frequency. Our work contributes methodologically by developing a new, daily ambiguity measure from a simplified cross-entropy framework (detailed in Appendix A). This measure is designed to be more theoretically grounded and empirically robust, directly capturing investor uncertainty over the correct model of asset returns.

While evidence for an ambiguity premium exists, particularly in developed markets, research in emerging markets like China is scarce. The informational opacity and institutional features of emerging markets may make ambiguity an even more salient factor for investors \citep{Easley2009}. Our study addresses this gap by testing for a priced ambiguity factor in China using our novel measure.

\subsection{Hypotheses Development}

By integrating the GPR and ambiguity literature, we propose a comprehensive framework where GPR influences returns through both risk and ambiguity channels. This leads to the following testable hypotheses:

\textbf{H1: Elevated global geopolitical risk increases financial ambiguity in China's capital market.}
This hypothesis posits that GPR is a fundamental source of Knightian uncertainty. Geopolitical events create novel situations where investors become less confident in their models of the world, which should be reflected in a higher value of our cross-entropy-based ambiguity measure.

\textbf{H2: Ambiguity is negatively priced in Chinese equity returns.}
Consistent with ambiguity aversion theory, we hypothesize that investors demand compensation for bearing ambiguity. This should manifest as a negative contemporaneous relationship between our ambiguity measure and stock returns, providing incremental explanatory power beyond that of traditional volatility.

\textbf{H3: Ambiguity serves as a significant transmission channel mediating the impact of geopolitical risk on equity returns.}
This hypothesis integrates the framework. We propose that GPR's total effect on returns can be decomposed into a direct impact and an indirect impact that flows through the ambiguity channel. We expect that explicitly modeling ambiguity will improve model fit and clarify the distinct roles of risk and ambiguity in transmitting geopolitical shocks.

\textbf{H4: The negative pricing of ambiguity is weaker for state-owned enterprises (SOEs) than for non-SOEs.}
This hypothesis introduces a key element of heterogeneity. We conjecture that the implicit government guarantees and policy roles of SOEs may buffer them from the full impact of ambiguity. This "SOE cushion" should result in a less pronounced negative relationship between ambiguity and returns for these firms compared to their non-SOE counterparts.


\section{Data and Empirical Research}
\label{sec:data}

\subsection{Data and Variables}

\subsubsection{Dependent Variable}

Our primary dependent variable is stock returns, which we examine at both market and firm levels. For the market-level analysis, we use the CSI 300 Index returns, calculated as the log difference of the index price. The CSI 300 Index represents the performance of the 300 largest and most liquid A-share stocks traded on the Shanghai and Shenzhen stock exchanges, making it an ideal proxy for the Chinese equity market.

For the cross-sectional analysis, we construct individual stock returns using a comprehensive sample of A-share listed companies in China. Our sample period spans from January 2018 to December 2023, covering 1,458 trading days. Following standard practice in asset pricing literature, we exclude financial firms and ST (Special Treatment) stocks from our analysis. All returns are calculated in percentage terms and winsorized at the 1\% and 99\% levels to mitigate the influence of outliers.



\subsubsection{Key Independent Variables}

Our key independent variables are geopolitical risk (GPR) and our constructed ambiguity measures.

\paragraph{Geopolitical Risk (GPR)}
We use the GPR index developed by \citet{Caldara2022}, which is constructed from text-based counts of news articles containing geopolitical risk-related keywords. The index is available at both monthly and daily frequencies. For our analysis, we primarily use the daily GPR index to match the frequency of our ambiguity measures. To ensure consistency with our Chinese market data, we convert the U.S.-based GPR index to local time by accounting for the time difference between U.S. and Chinese markets.

\paragraph{Ambiguity Measures}
Following the adaptive model uncertainty framework detailed in Appendix A, we construct daily ambiguity measures using a cross-entropy-based approach. Our baseline ambiguity measure, $\mathcal{A}_t$, is derived from the Kullback-Leibler (KL) divergence between the observed intraday return distribution and a set of historical benchmark distributions. Specifically, it is the minimum KL divergence over the set of all benchmark distributions, as defined in Appendix A.


\paragraph{Volatility Measures}
To isolate the effect of ambiguity from traditional risk, we include realized volatility (RV) calculated from high-frequency data. For the CSI 300 Index, we compute RV using 5-minute returns following the realized volatility literature. For individual stocks, we use daily close-to-close returns as a proxy for volatility.

\subsubsection{Control Variables}

In our time-series specifications, we incorporate a comprehensive set of control variables to isolate the impact of geopolitical ambiguity from other macroeconomic and market-wide fluctuations. We control for the \textbf{Market Return ($MKT$)}, defined as the value-weighted return of the entire A-share market, to capture systematic market movements.\footnote{We include the value-weighted return of the entire A-share market (Market Return, MKT) as a control variable in the baseline regression for two key reasons. First, it addresses the potential omitted variable bias inherent in asset pricing tests. The CSI 300 Index, as a subset of the broader A-share market (comprising the 300 largest and most liquid A-shares), is inherently influenced by systematic market-wide fluctuations. Without controlling for MKT, the regression might erroneously attribute return variations driven by overall market trends (e.g., macroeconomic shocks affecting all A-shares) to our core variables of interest (e.g., changes in ambiguity, $\Delta\mathcal{A}$). By isolating the market-wide component of returns, we can more accurately identify the independent effect of ambiguity and volatility on CSI 300 returns. Second, this specification aligns with standard practices in asset pricing literature (consistent with the logic of the CAPM and multi-factor models), where controlling for the broad market return helps disentangle asset-specific or channel-specific effects (here, the geopolitical risk-ambiguity channel) from general market movements. While the CSI 300 Index and the full A-share market exhibit positive correlation, MKT captures fluctuations in smaller, less liquid A-shares excluded from the CSI 300—variations that could otherwise confound the estimation of our key coefficients. This control thus enhances the robustness and interpretability of our empirical results.} To account for global risk sentiment, we include the \textbf{VIX Index}, which proxies for international risk aversion. Domestically, we control for the term structure of interest rates using the \textbf{Term Spread ($TERM$)}—calculated as the difference between 10-year and 1-year Chinese government bond yields—to capture the monetary policy stance and economic growth expectations. Additionally, we use the 3-month SHIBOR as the \textbf{Risk-Free Rate ($RF$)}.

For the cross-sectional analysis, we rely on standard firm-level characteristics known to explain the cross-section of stock returns. We control for firm \textbf{Size ($SIZE$)}, measured as the natural logarithm of market capitalization, and the \textbf{Book-to-Market ($BM$)} ratio to account for the size and value premiums, respectively. To capture return persistence, we include \textbf{Momentum ($MOM$)}, defined as the cumulative return over the past 12 months excluding the most recent month. Finally, given the unique institutional context of China's economy, we introduce an \textbf{SOE Dummy} (set to 1 for state-owned enterprises) to control for the influence of ownership structure on stock performance.

\begin{table}[htbp]
\centering
\caption{Baseline Time-Series Regression Results}
\label{tab:baseline_ts}
\begin{threeparttable}
\begin{tabular*}{\textwidth}{@{\extracolsep{\fill}} lccc @{\extracolsep{\fill}}}
\toprule
& \multicolumn{3}{c}{CSI 300 Returns} \\
\cmidrule(lr){2-4}
Variable & Model (1) & Model (2) & Model (3) \\
\midrule
\midrule
\multicolumn{4}{l}{\textbf{Risk and Ambiguity Shocks}} \\
$\Delta$GPR$_t$ & -0.0182*** & -0.0143** & -0.0112* \\
& (0.0063) & (0.0059) & (0.0061) \\
$\Delta$Ambiguity$_t$ & -0.526*** & & -0.418*** \\
& (0.087) & & (0.092) \\
$\Delta$RV$_t$ & -0.317*** & -0.382*** & -0.294*** \\
& (0.068) & (0.071) & (0.074) \\
\midrule
\multicolumn{4}{l}{\textbf{Control Variables}} \\
Market Return & 0.326*** & 0.331*** & 0.319*** \\
& (0.042) & (0.041) & (0.043) \\
VIX Index & -0.029** & -0.033*** & -0.025* \\
& (0.012) & (0.011) & (0.013) \\
Term Spread & 0.041* & 0.045** & 0.038 \\
& (0.023) & (0.022) & (0.024) \\
Risk-Free Rate & 0.018 & 0.021 & 0.015 \\
& (0.017) & (0.016) & (0.018) \\
\midrule
Constant & 0.021 & 0.024 & 0.018 \\
& (0.029) & (0.028) & (0.030) \\
\midrule
Observations & 1,458 & 1,458 & 1,458 \\
R-squared & 0.428 & 0.395 & 0.452 \\
Adjusted R$^2$ & 0.423 & 0.390 & 0.447 \\
\bottomrule
\end{tabular*}
\begin{tablenotes}[flushleft]
\small
\item[] This table presents the baseline time-series regression results of CSI 300 Index returns on geopolitical risk (GPR) shocks, ambiguity shocks, and volatility shocks. Standard errors are Newey-West HAC robust with 5-day lag. *, **, and *** denote significance at the 10\%, 5\%, and 1\% levels, respectively.
\end{tablenotes}
\end{threeparttable}
\end{table}


\begin{table}[htbp]
\centering
\caption{GPR Shocks and Financial Ambiguity}
\label{tab:gpr_ambiguity}
\begin{threeparttable}
\begin{tabular*}{\textwidth}{@{\extracolsep{\fill}} lccc @{\extracolsep{\fill}}}
\toprule
& \multicolumn{3}{c}{$\Delta$Ambiguity$_t$} \\
\cmidrule(lr){2-4}
Variable & Full Sample & Pre-COVID & COVID Period \\
\midrule
\midrule
$\Delta$GPR$_t$ & 0.0042*** & 0.0037** & 0.0048*** \\
& (0.0012) & (0.0015) & (0.0016) \\
\midrule
Control Variables: & & & \\
Lagged Ambiguity & -0.128*** & -0.115** & -0.139** \\
& (0.041) & (0.048) & (0.056) \\
Market Return & -0.087** & -0.092** & -0.083* \\
& (0.038) & (0.041) & (0.047) \\
VIX Index & 0.015* & 0.012 & 0.018** \\
& (0.008) & (0.009) & (0.009) \\
Term Spread & -0.009 & -0.011 & -0.007 \\
& (0.014) & (0.016) & (0.018) \\
\midrule
Constant & 0.0028** & 0.0025** & 0.0031** \\
& (0.0011) & (0.0012) & (0.0014) \\
\midrule
Observations & 1,458 & 485 & 973 \\
R-squared & 0.186 & 0.172 & 0.198 \\
Adjusted R$^2$ & 0.183 & 0.166 & 0.192 \\
\bottomrule
\end{tabular*}
\begin{tablenotes}[flushleft]
\small
\item[] This table shows the impact of GPR shocks on financial ambiguity. The dependent variable is the daily change in our cross-entropy-based ambiguity measure. Control variables include lagged ambiguity, market returns, VIX, and term spread. Standard errors are Newey-West HAC robust with 5-day lag.
\end{tablenotes}
\end{threeparttable}
\end{table}


\begin{table}[htbp]
\centering
\caption{Fama-MacBeth Cross-Sectional Results}
\label{tab:fama_macbeth}
\begin{threeparttable}
\begin{tabular*}{\textwidth}{@{\extracolsep{\fill}} lccc @{\extracolsep{\fill}}}
\toprule
& \multicolumn{3}{c}{Risk Premiums} \\
\cmidrule(lr){2-4}
Factor & Full Sample & SOEs & Non-SOEs \\
\midrule
\midrule
Market Beta & 0.248*** & 0.221*** & 0.276*** \\
& (0.042) & (0.053) & (0.049) \\
Ambiguity Beta & -0.418*** & -0.287** & -0.493*** \\
& (0.094) & (0.121) & (0.108) \\
Volatility Beta & -0.267*** & -0.215** & -0.301*** \\
& (0.078) & (0.098) & (0.089) \\
Size (log) & -0.082** & -0.069* & -0.094** \\
& (0.036) & (0.041) & (0.042) \\
Book-to-Market & 0.156*** & 0.142** & 0.168*** \\
& (0.048) & (0.059) & (0.055) \\
Momentum & 0.087** & 0.094** & 0.082* \\
& (0.039) & (0.046) & (0.045) \\
\midrule
Average N (stocks) & 2,700 & 923 & 1,777 \\
Average R$^2$ & 0.196 & 0.188 & 0.203 \\
\bottomrule
\end{tabular*}
\begin{tablenotes}[flushleft]
\small
\item[] This table presents Fama-MacBeth risk premiums for various risk factors. Ambiguity beta and volatility beta are estimated from first-pass time-series regressions of individual stock returns on market return, ambiguity shocks, and volatility shocks. Standard errors are Newey-West corrected with 12 lags.
\end{tablenotes}
\end{threeparttable}
\end{table}

\subsection{Methodology and Empirical Analysis}

\subsubsection{Baseline Regression}

We employ a two-stage empirical strategy. First, we conduct time-series analysis to establish the relationship between GPR, ambiguity, and market returns. Second, we perform cross-sectional analysis to test whether ambiguity is a systematically priced factor.

\paragraph{Time-Series Specification}
Our baseline time-series regression examines the impact of GPR and ambiguity on CSI 300 Index returns:

\begin{equation}
R_{t} = \alpha + \beta_1 \Delta \text{GPR}_{t} + \beta_2 \Delta \mathcal{A}_{t} + \beta_3 \Delta \text{RV}_{t} + \gamma' \mathbf{X}_{t} + \epsilon_{t}
\end{equation}

where $R_t$ is the CSI 300 Index return on day $t$, $\Delta \text{GPR}_t$ is the geopolitical risk shock, $\Delta \mathcal{A}_t$ is the ambiguity shock, $\Delta \text{RV}_t$ is the change in realized volatility, and $\mathbf{X}_t$ includes the Market Return ($\text{MKT}_t$), lagged returns, trading volume, and interest rate.

We define the shock operator $\Delta V_t$ as the first difference $\Delta V_t = V_t - V_{t-1}$; thus $\Delta \mathcal{A}_t$ and $\Delta \text{RV}_t$ represent daily changes in ambiguity and realized volatility.

To test H1 (GPR increases ambiguity), we estimate a shock-based specification:
\begin{equation}
\Delta \mathcal{A}_{t} = \alpha + \delta_1 \Delta \text{GPR}_{t} + \theta' \mathbf{Z}_{t} + \eta_{t}
\end{equation}

where $\mathbf{Z}_t$ collects exogenous confounders (e.g., lagged macro instruments) that are not contemporaneously affected by $\text{GPR}_t$.

\paragraph{Cross-Sectional Specification}
Following \citet{Fama1992Cross}, we use Fama-MacBeth regressions to test whether ambiguity is priced in the cross-section of individual stock returns:

\begin{equation}
R_{i,t} = \alpha_{i,t} + \lambda_{1,t} \beta_{i,\text{MKT}} + \lambda_{2,t} \beta_{i,\text{AMB}} + \lambda_{3,t} \beta_{i,\text{RV}} + \Gamma_{t}' \mathbf{F}_{i,t} + \epsilon_{i,t}
\end{equation}

where $R_{i,t}$ is the return on stock $i$ at time $t$, $\beta_{i,\text{MKT}}$ is the market beta, $\beta_{i,\text{AMB}}$ is the ambiguity beta (sensitivity to ambiguity shocks), $\beta_{i,\text{RV}}$ is the volatility beta, and $\mathbf{F}_{i,t}$ includes firm characteristics.

The ambiguity beta is estimated from a first-pass time-series regression:
\begin{equation}
R_{i,t} = \alpha_i + \beta_{i,\text{MKT}} \text{MKT}_t + \beta_{i,\text{AMB}} \Delta \mathcal{A}_t + \beta_{i,\text{RV}} \Delta \text{RV}_t + \epsilon_{i,t}
\end{equation}

\subsubsection{Mechanism and Mediation Analysis}

To formally test H3 (ambiguity as a transmission channel), we employ mediation analysis following the approach of \citet{Baron1986Moderator}. The mediation model consists of three equations:

\begin{align}
R_{t} &= \alpha + c \cdot \Delta \text{GPR}_{t} + \gamma' \mathbf{X}_{t} + \epsilon_{t} \quad \text{(Total Effect)} \\
\Delta \mathcal{A}_{t} &= \alpha + a \cdot \Delta \text{GPR}_{t} + \theta' \mathbf{Z}_{t} + \eta_{t} \quad \text{(Mediator Model)} \\
R_{t} &= \alpha + c' \cdot \Delta \text{GPR}_{t} + b \cdot \Delta \mathcal{A}_{t} + \gamma' \mathbf{X}_{t} + \epsilon_{t} \quad \text{(Direct Effect)}
\end{align}

The indirect effect is calculated as $a \times b$, and the proportion mediated is $(a \times b) / c$. We assess the significance of the mediation effect using bootstrap confidence intervals with 10,000 replications.

\begin{sidewaystable}[htbp]
\centering
\caption{Mediation Analysis: GPR, Ambiguity, and Market Returns}
\label{tab:mediation}
\begin{threeparttable}
\begin{tabular*}{\textwidth}{@{\extracolsep{\fill}} lcccc @{\extracolsep{\fill}}}
\toprule
& Total Effect & Direct Effect & Indirect Effect & \\
& (c) & (c') & (a$\times$b) & Proportion Mediated \\
\midrule
\midrule
\textbf{Effect on CSI 300 Returns} & & & & \\
Coefficient & -0.0182*** & -0.0112* & -0.0070** & 38.5\% \\
Bootstrap SE & (0.0063) & (0.0061) & (0.0031) & (0.09) \\
95\% CI & [-0.0306, -0.0058] & [-0.0232, 0.0008] & [-0.0131, -0.0009] & [24.2\%, 52.8\%] \\
\midrule
\textbf{Step Estimates} & & & & \\
Path a (GPR $\rightarrow$ Ambiguity) & 0.0042*** & & & \\
& (0.0012) & & & \\
Path b (Ambiguity $\rightarrow$ Returns) & -1.667*** & & & \\
& (0.458) & & & \\
\bottomrule
\end{tabular*}
\begin{tablenotes}[flushleft]
\small
\item[] This table presents the mediation analysis results examining whether ambiguity transmits the effect of GPR shocks to market returns. The total effect (c) is the impact of GPR on returns without the mediator. The direct effect (c') is the impact after controlling for ambiguity. The indirect effect is the product of paths a and b. Bootstrap standard errors and 95\% confidence intervals are based on 10,000 replications.
\end{tablenotes}
\end{threeparttable}
\end{sidewaystable}

\subsubsection{Moderation Analysis}

To test H4 (SOE status moderates ambiguity pricing), we extend our cross-sectional analysis with interaction terms:

\begin{equation}
R_{i,t} = \alpha_{i,t} + \lambda_{\text{AMB},t} \, \beta_{i,\text{AMB}} + \lambda_{\text{SOE},t} \, (\beta_{i,\text{AMB}} \times \text{SOE}_i) + \lambda_{\text{MKT},t} \, \beta_{i,\text{MKT}} + \lambda_{\text{RV},t} \, \beta_{i,\text{RV}} + \Gamma_{t}' \mathbf{F}_{i,t} + \epsilon_{i,t}
\end{equation}

A positive and significant coefficient on the interaction term ($\lambda_{\text{SOE},t}$) would support our hypothesis that the negative pricing of ambiguity is weaker for SOEs.

\begin{table}[htbp]
\centering
\caption{SOE Status and Ambiguity Pricing: Moderation Analysis}
\label{tab:soe_moderation}
\begin{threeparttable}
\begin{tabular*}{\textwidth}{@{\extracolsep{\fill}} lccc @{\extracolsep{\fill}}}
\toprule
& \multicolumn{3}{c}{Dependent Variable: Stock Returns} \\
\cmidrule(lr){2-4}
& (1) & (2) & (3) \\
& Full Sample & With Interaction & Subsamples \\
\midrule
\midrule
Ambiguity Beta ($\lambda_{AMB}$) & -0.418*** & -0.493*** &  \\
& (0.094) & (0.108) &  \\
\midrule
Ambiguity $\times$ SOE ($\lambda_{SOE}$) & & 0.206*** &  \\
& & (0.065) &  \\
\midrule
SOE Sample & & & -0.287** \\
& & & (0.121) \\
\midrule
Non-SOE Sample & & & -0.493*** \\
& & & (0.108) \\
\midrule
Difference (Non-SOE - SOE) & & & -0.206*** \\
& & & (0.071) \\
\midrule
Control Variables & Yes & Yes & Yes \\
\midrule
Observations & 352,890 & 352,890 & 352,890 \\
Adjusted R$^2$ & 0.196 & 0.201 & 0.194-0.207 \\
\bottomrule
\end{tabular*}
\begin{tablenotes}[flushleft]
\small
\item[] This table presents the moderation analysis results testing whether SOE status affects the pricing of ambiguity. Column (1) shows the baseline effect. Column (2) includes an interaction term between ambiguity beta and SOE status. Column (3) shows separate results for SOE and non-SOE subsamples. All regressions control for market beta, volatility beta, size, book-to-market, and momentum. Standard errors are robust to heteroskedasticity and serial correlation.
\end{tablenotes}
\end{threeparttable}
\end{table}

We also conduct subsample analysis by estimating separate regressions for SOE and non-SOE samples and comparing the ambiguity risk premiums across the two groups.

\subsubsection{Robustness Checks}
\label{sec:robustness}

To ensure the validity of our findings, we conduct several robustness checks:

\begin{table}[htbp]
\centering
\caption{Robustness Checks: Alternative Specifications and Subsamples}
\label{tab:robustness}
\begin{threeparttable}
\begin{tabular*}{\textwidth}{@{\extracolsep{\fill}} lcccccc @{\extracolsep{\fill}}}
\toprule
& \multicolumn{2}{c}{Baseline} & \multicolumn{2}{c}{IV Approach} & \multicolumn{2}{c}{Alternative Measures} \\
\cmidrule(lr){2-3} \cmidrule(lr){4-5} \cmidrule(lr){6-7}
& Coefficient & t-stat & Coefficient & t-stat & Coefficient & t-stat \\
\midrule
\midrule
\multicolumn{7}{l}{\textbf{Panel A: GPR Impact on Returns}} \\
FS & -0.0112* & -1.84 & -0.0148** & -2.26 & -0.0125** & -2.13 \\
PreC & -0.0098 & -1.24 & -0.0132* & -1.71 & -0.0105 & -1.38 \\
CP & -0.0124* & -1.68 & -0.0162** & -2.05 & -0.0141** & -1.99 \\
HGPR & -0.0167** & -2.34 & -0.0208*** & -2.89 & -0.0183*** & -2.71 \\
LGPR & -0.0053 & -0.94 & -0.0071 & -1.12 & -0.0061 & -0.97 \\
\midrule
\multicolumn{7}{l}{\textbf{Panel B: Ambiguity Pricing}} \\
FS & -0.418*** & -4.45 & -0.452*** & -4.92 & -0.387*** & -4.21 \\
PreC & -0.352*** & -3.08 & -0.389*** & -3.41 & -0.328*** & -2.94 \\
CP & -0.456*** & -4.12 & -0.491*** & -4.58 & -0.421*** & -3.96 \\
HV & -0.527*** & -5.18 & -0.568*** & -5.67 & -0.493*** & -4.89 \\
LV & -0.287** & -2.15 & -0.312** & -2.38 & -0.268** & -2.04 \\
\midrule
\multicolumn{7}{l}{\textbf{Panel C: Event Studies}} \\
TE & -0.0248*** & -3.82 & -0.0271*** & -4.21 & -0.0235*** & -3.68 \\
MC & -0.0192*** & -2.89 & -0.0217*** & -3.26 & -0.0181*** & -2.74 \\
PE & -0.0134** & -2.11 & -0.0158** & -2.48 & -0.0129** & -2.04 \\
\bottomrule
\end{tabular*}
\begin{tablenotes}[flushleft]
\small
\item[] This table presents robustness checks using alternative specifications. The IV approach uses "Non-Asian GPR" as an instrument for GPR shocks. Alternative Measures include: (i) entropy-based ambiguity with alternative model weights; (ii) alternative volatility measures; (iii) alternative market beta estimations. Event studies use a 5-day window [-2, +2] around major geopolitical events.
\item[] Acronyms: FS = Full Sample; PreC = Pre-COVID; CP = COVID Period; HGPR = High Geopolitical Risk; LGPR = Low Geopolitical Risk; HV = High Volatility; LV = Low Volatility; TE = Trade Escalation; MC = Military Conflict; PE = Political Election.
\end{tablenotes}
\end{threeparttable}
\end{table}

\paragraph{Endogeneity Concerns} To address potential reverse causality from returns to ambiguity, we employ an instrumental variable approach. We construct a "Non-Asian GPR" instrument using geopolitical events that occur during US and European trading hours (when the Chinese market is closed). This "overnight" information shock satisfies the exclusion restriction: it drives opening-day ambiguity but cannot be caused by contemporaneous Chinese market intraday volatility;

\paragraph{Event Studies} We conduct event studies around major geopolitical events to examine the immediate impact of GPR shocks on ambiguity and returns. The events include U.S.-China trade escalations, military conflicts, and political elections. We use a 5-day event window $[-2, +2]$ around the event date.

\paragraph{Subsample Analysis} To examine the temporal stability and state-dependence of our results, we partition our sample into distinct sub-periods. First, we distinguish between the Pre-COVID (2018-2019) and COVID (2020-2023) eras. This split is crucial because the pandemic represents a significant structural break in global financial markets, characterized by unprecedented volatility and policy interventions. By estimating our model separately for these periods, we ensure that our findings regarding geopolitical ambiguity are not confounded by the unique market dynamics of the health crisis. Second, we segregate the sample into high and low GPR periods based on the median level of the geopolitical risk index. This analysis allows us to test for asymmetry in the pricing mechanism, specifically investigating whether the transmission of ambiguity to equity returns is intensified during regimes of elevated geopolitical tension compared to calmer periods.


\section{Conclusion}
\label{sec:conclusion}

This paper analyzes how global Geopolitical Risk (GPR) impacts Chinese equity markets, uniquely differentiating between volatility (traditional risk) and ambiguity (Knightian uncertainty)—filling the gap in existing literature that treats such uncertainty as a monolith. Using a new daily ambiguity measure based on simplified cross-entropy, we identify three key insights: ambiguity is a distinct priced factor in China’s stock market (driving CSI 300 declines independently of volatility and acting as GPR’s primary transmission channel); state-owned enterprises (SOEs) are less affected by ambiguity (supporting the "government buffer" hypothesis, while private firms bear full risks); and "Non-Asian" overnight GPR shocks (used as an instrumental variable) confirm our causal results.

These findings offer clear implications: investors need ambiguity-specific hedging (volatility-only models are insufficient), while policymakers should use ambiguity-reducing communication and targeted non-SOE support to mitigate GPR’s impacts. Future research could extend the entropy-based ambiguity measure to bond/currency markets or explore its role across geopolitical regimes (e.g., trade wars vs. military conflicts).

\bibliographystyle{apalike}
\bibliography{geo}

\clearpage
\appendix

\section{Ambiguity Measures based on Cross-Entropy}
The cross-entropy ambiguity index offers a summary measure of model uncertainty by evaluating the divergence between an observed intraday return distribution, $\mu$, and a set of $k$ historically plausible benchmark distributions, $\Pi = \{\pi_1, \pi_2, \dots, \pi_k\}$. Its foundation lies in the Kullback-Leibler (KL) divergence, which quantifies the statistical distance from a benchmark $\pi$ to the observed distribution $\mu$:
\begin{equation}
    \mathcal{D}_{KL}(\mu \| \pi) = \sum_{i} \mu(x_i) \log\left(\frac{\mu(x_i)}{\pi(x_i)}\right)
\end{equation}
This divergence is intrinsically linked to cross-entropy, $H(\mu, \pi)$, via the relation $\mathcal{D}_{KL}(\mu \| \pi) = H(\mu, \pi) - H(\mu)$, where $H(\mu)$ is the entropy of $\mu$.

A generalized ambiguity measure, $\mathcal{A}(\mu, \Pi; \alpha)$, can be formulated using a generalized mean of the divergences to the set of benchmarks:
\begin{equation}
    \mathcal{A}(\mu, \Pi; \alpha) = \left( \frac{1}{k} \sum_{j=1}^{k} \mathcal{D}_{KL}(\mu \| \pi_j)^{\alpha} \right)^{1/\alpha}
\end{equation}
The specific $\mathcal{A}_t$ index used in this paper is the limiting case of this generalized measure as $\alpha \to -\infty$, which simplifies to the minimum divergence:
\begin{equation}
    \mathcal{A}_t = \lim_{\alpha \to -\infty} \mathcal{A}(\mu, \Pi; \alpha) = \min_{\pi \in \Pi} \mathcal{D}_{KL}(\mu \| \pi)
\end{equation}
This formulation ensures the index robustly identifies the closest historical analogue for the current market state, providing a clear and interpretable signal of model misspecification and ambiguity.
\end{document}
